\section{Matriz de la prueba de la PLJ}\label{sec:matriz}

Estado final de la matriz de la prueba de la PLJ para $\rho' = \rho \cup \{\text{Llave}\}$. Por espacio se presenta traspuesta: cada fila es un atributo de $U$ y cada columna un esquema. Se omiten los subíndices: en la fila del atributo $A_j$, \emph{a} es $a_j$ y \emph{b} es el $b_{ij}$ de la columna. Una \emph{a} normal indica que el atributo pertenece al esquema ($A_j \in U_i$), así que la celda tiene $a_j$ desde el inicio del algoritmo, como se describe en la sección~\ref{sec:plj}. Una \textbf{\color{azuldos}a} resaltada indica que el atributo no pertenece al esquema: la celda empezó con $b_{ij}$ y el algoritmo la cambió por $a_j$. Una {\color{grisclaro}b} no cambió a $a_j$, aunque pudo igualarse a otra $b$. La columna Llave es la única formada solo por símbolos $a_j$.

% Generado por assets/source/verificacion.py. No editar a mano.
{\scriptsize\setlength{\tabcolsep}{2.2pt}\renewcommand{\arraystretch}{1.05}\filas
\begin{longtable}{lccccccccccccccccccccc}
\rowcolor{azul}\cab{Atributo} & \rotatebox{90}{\cab{Programa}} & \rotatebox{90}{\cab{Asignatura}} & \rotatebox{90}{\cab{Coordinador}} & \rotatebox{90}{\cab{Autorización}} & \rotatebox{90}{\cab{Estudiante}} & \rotatebox{90}{\cab{Matrícula}} & \rotatebox{90}{\cab{Jefe}} & \rotatebox{90}{\cab{Bloque}} & \rotatebox{90}{\cab{Semestre}} & \rotatebox{90}{\cab{Horario}} & \rotatebox{90}{\cab{Validación}} & \rotatebox{90}{\cab{Revisión}} & \rotatebox{90}{\cab{Sustitución}} & \rotatebox{90}{\cab{Aula}} & \rotatebox{90}{\cab{Franja}} & \rotatebox{90}{\cab{Sesión}} & \rotatebox{90}{\cab{Asignación}} & \rotatebox{90}{\cab{Fecha}} & \rotatebox{90}{\cab{Asistencia}} & \rotatebox{90}{\cab{Solicitud}} & \rotatebox{90}{\cab{Llave}} \\
\endhead
\#P & a & a & \color{grisclaro}b & a & \color{grisclaro}b & a & \color{grisclaro}b & \textbf{\color{azuldos}a} & \color{grisclaro}b & a & \textbf{\color{azuldos}a} & \textbf{\color{azuldos}a} & \textbf{\color{azuldos}a} & \color{grisclaro}b & \color{grisclaro}b & \textbf{\color{azuldos}a} & \textbf{\color{azuldos}a} & \color{grisclaro}b & \textbf{\color{azuldos}a} & \color{grisclaro}b & \textbf{\color{azuldos}a} \\
NombreP & a & \textbf{\color{azuldos}a} & \color{grisclaro}b & \textbf{\color{azuldos}a} & \color{grisclaro}b & \textbf{\color{azuldos}a} & \color{grisclaro}b & \textbf{\color{azuldos}a} & \color{grisclaro}b & \textbf{\color{azuldos}a} & \textbf{\color{azuldos}a} & \textbf{\color{azuldos}a} & \textbf{\color{azuldos}a} & \color{grisclaro}b & \color{grisclaro}b & \textbf{\color{azuldos}a} & \textbf{\color{azuldos}a} & \color{grisclaro}b & \textbf{\color{azuldos}a} & \color{grisclaro}b & \textbf{\color{azuldos}a} \\
Plan & a & \textbf{\color{azuldos}a} & \color{grisclaro}b & \textbf{\color{azuldos}a} & \color{grisclaro}b & \textbf{\color{azuldos}a} & \color{grisclaro}b & \textbf{\color{azuldos}a} & \color{grisclaro}b & \textbf{\color{azuldos}a} & \textbf{\color{azuldos}a} & \textbf{\color{azuldos}a} & \textbf{\color{azuldos}a} & \color{grisclaro}b & \color{grisclaro}b & \textbf{\color{azuldos}a} & \textbf{\color{azuldos}a} & \color{grisclaro}b & \textbf{\color{azuldos}a} & \color{grisclaro}b & \textbf{\color{azuldos}a} \\
\#A & \color{grisclaro}b & a & \color{grisclaro}b & \color{grisclaro}b & \color{grisclaro}b & \color{grisclaro}b & \color{grisclaro}b & a & \color{grisclaro}b & \color{grisclaro}b & \color{grisclaro}b & \color{grisclaro}b & \color{grisclaro}b & \color{grisclaro}b & \color{grisclaro}b & \textbf{\color{azuldos}a} & \color{grisclaro}b & \color{grisclaro}b & \textbf{\color{azuldos}a} & \color{grisclaro}b & \textbf{\color{azuldos}a} \\
NombreA & \color{grisclaro}b & a & \color{grisclaro}b & \color{grisclaro}b & \color{grisclaro}b & \color{grisclaro}b & \color{grisclaro}b & \textbf{\color{azuldos}a} & \color{grisclaro}b & \color{grisclaro}b & \color{grisclaro}b & \color{grisclaro}b & \color{grisclaro}b & \color{grisclaro}b & \color{grisclaro}b & \textbf{\color{azuldos}a} & \color{grisclaro}b & \color{grisclaro}b & \textbf{\color{azuldos}a} & \color{grisclaro}b & \textbf{\color{azuldos}a} \\
\#C & \color{grisclaro}b & \color{grisclaro}b & a & a & \color{grisclaro}b & \color{grisclaro}b & \color{grisclaro}b & \color{grisclaro}b & \color{grisclaro}b & \color{grisclaro}b & a & \color{grisclaro}b & \color{grisclaro}b & \color{grisclaro}b & \color{grisclaro}b & \color{grisclaro}b & \color{grisclaro}b & \color{grisclaro}b & \color{grisclaro}b & \color{grisclaro}b & a \\
NombreC & \color{grisclaro}b & \color{grisclaro}b & a & \textbf{\color{azuldos}a} & \color{grisclaro}b & \color{grisclaro}b & \color{grisclaro}b & \color{grisclaro}b & \color{grisclaro}b & \color{grisclaro}b & \textbf{\color{azuldos}a} & \color{grisclaro}b & \color{grisclaro}b & \color{grisclaro}b & \color{grisclaro}b & \color{grisclaro}b & \color{grisclaro}b & \color{grisclaro}b & \color{grisclaro}b & \color{grisclaro}b & \textbf{\color{azuldos}a} \\
Especialidad & \color{grisclaro}b & \color{grisclaro}b & a & \textbf{\color{azuldos}a} & \color{grisclaro}b & \color{grisclaro}b & \color{grisclaro}b & \color{grisclaro}b & \color{grisclaro}b & \color{grisclaro}b & \textbf{\color{azuldos}a} & \color{grisclaro}b & \color{grisclaro}b & \color{grisclaro}b & \color{grisclaro}b & \color{grisclaro}b & \color{grisclaro}b & \color{grisclaro}b & \color{grisclaro}b & \color{grisclaro}b & \textbf{\color{azuldos}a} \\
PGenerar & \color{grisclaro}b & \color{grisclaro}b & \color{grisclaro}b & a & \color{grisclaro}b & \color{grisclaro}b & \color{grisclaro}b & \color{grisclaro}b & \color{grisclaro}b & \color{grisclaro}b & \textbf{\color{azuldos}a} & \color{grisclaro}b & \color{grisclaro}b & \color{grisclaro}b & \color{grisclaro}b & \color{grisclaro}b & \color{grisclaro}b & \color{grisclaro}b & \color{grisclaro}b & \color{grisclaro}b & \textbf{\color{azuldos}a} \\
PValidar & \color{grisclaro}b & \color{grisclaro}b & \color{grisclaro}b & a & \color{grisclaro}b & \color{grisclaro}b & \color{grisclaro}b & \color{grisclaro}b & \color{grisclaro}b & \color{grisclaro}b & \textbf{\color{azuldos}a} & \color{grisclaro}b & \color{grisclaro}b & \color{grisclaro}b & \color{grisclaro}b & \color{grisclaro}b & \color{grisclaro}b & \color{grisclaro}b & \color{grisclaro}b & \color{grisclaro}b & \textbf{\color{azuldos}a} \\
\#E & \color{grisclaro}b & \color{grisclaro}b & \color{grisclaro}b & \color{grisclaro}b & a & a & \color{grisclaro}b & \color{grisclaro}b & \color{grisclaro}b & \color{grisclaro}b & \color{grisclaro}b & \color{grisclaro}b & \color{grisclaro}b & \color{grisclaro}b & \color{grisclaro}b & \color{grisclaro}b & a & \color{grisclaro}b & a & a & \textbf{\color{azuldos}a} \\
NombreE & \color{grisclaro}b & \color{grisclaro}b & \color{grisclaro}b & \color{grisclaro}b & a & \textbf{\color{azuldos}a} & \color{grisclaro}b & \color{grisclaro}b & \color{grisclaro}b & \color{grisclaro}b & \color{grisclaro}b & \color{grisclaro}b & \color{grisclaro}b & \color{grisclaro}b & \color{grisclaro}b & \color{grisclaro}b & \textbf{\color{azuldos}a} & \color{grisclaro}b & \textbf{\color{azuldos}a} & \textbf{\color{azuldos}a} & \textbf{\color{azuldos}a} \\
Edad & \color{grisclaro}b & \color{grisclaro}b & \color{grisclaro}b & \color{grisclaro}b & a & \textbf{\color{azuldos}a} & \color{grisclaro}b & \color{grisclaro}b & \color{grisclaro}b & \color{grisclaro}b & \color{grisclaro}b & \color{grisclaro}b & \color{grisclaro}b & \color{grisclaro}b & \color{grisclaro}b & \color{grisclaro}b & \textbf{\color{azuldos}a} & \color{grisclaro}b & \textbf{\color{azuldos}a} & \textbf{\color{azuldos}a} & \textbf{\color{azuldos}a} \\
Grupo & \color{grisclaro}b & \color{grisclaro}b & \color{grisclaro}b & \color{grisclaro}b & a & \textbf{\color{azuldos}a} & \color{grisclaro}b & \color{grisclaro}b & \color{grisclaro}b & \color{grisclaro}b & \color{grisclaro}b & \color{grisclaro}b & \color{grisclaro}b & \color{grisclaro}b & \color{grisclaro}b & \color{grisclaro}b & \textbf{\color{azuldos}a} & \color{grisclaro}b & \textbf{\color{azuldos}a} & \textbf{\color{azuldos}a} & \textbf{\color{azuldos}a} \\
\#J & \color{grisclaro}b & \color{grisclaro}b & \color{grisclaro}b & \color{grisclaro}b & \color{grisclaro}b & \color{grisclaro}b & a & \color{grisclaro}b & \color{grisclaro}b & \textbf{\color{azuldos}a} & \textbf{\color{azuldos}a} & a & \textbf{\color{azuldos}a} & \color{grisclaro}b & \color{grisclaro}b & \textbf{\color{azuldos}a} & \textbf{\color{azuldos}a} & \color{grisclaro}b & \textbf{\color{azuldos}a} & \color{grisclaro}b & \textbf{\color{azuldos}a} \\
NombreJ & \color{grisclaro}b & \color{grisclaro}b & \color{grisclaro}b & \color{grisclaro}b & \color{grisclaro}b & \color{grisclaro}b & a & \color{grisclaro}b & \color{grisclaro}b & \textbf{\color{azuldos}a} & \textbf{\color{azuldos}a} & \textbf{\color{azuldos}a} & \textbf{\color{azuldos}a} & \color{grisclaro}b & \color{grisclaro}b & \textbf{\color{azuldos}a} & \textbf{\color{azuldos}a} & \color{grisclaro}b & \textbf{\color{azuldos}a} & \color{grisclaro}b & \textbf{\color{azuldos}a} \\
\#B & \color{grisclaro}b & \color{grisclaro}b & \color{grisclaro}b & \color{grisclaro}b & \color{grisclaro}b & \color{grisclaro}b & \color{grisclaro}b & a & \color{grisclaro}b & \color{grisclaro}b & \color{grisclaro}b & \color{grisclaro}b & \color{grisclaro}b & \color{grisclaro}b & \color{grisclaro}b & a & \color{grisclaro}b & \color{grisclaro}b & a & \color{grisclaro}b & a \\
Descripcion & \color{grisclaro}b & \color{grisclaro}b & \color{grisclaro}b & \color{grisclaro}b & \color{grisclaro}b & \color{grisclaro}b & \color{grisclaro}b & a & \color{grisclaro}b & \color{grisclaro}b & \color{grisclaro}b & \color{grisclaro}b & \color{grisclaro}b & \color{grisclaro}b & \color{grisclaro}b & \textbf{\color{azuldos}a} & \color{grisclaro}b & \color{grisclaro}b & \textbf{\color{azuldos}a} & \color{grisclaro}b & \textbf{\color{azuldos}a} \\
Tipo & \color{grisclaro}b & \color{grisclaro}b & \color{grisclaro}b & \color{grisclaro}b & \color{grisclaro}b & \color{grisclaro}b & \color{grisclaro}b & a & \color{grisclaro}b & \color{grisclaro}b & \color{grisclaro}b & \color{grisclaro}b & \color{grisclaro}b & \color{grisclaro}b & \color{grisclaro}b & \textbf{\color{azuldos}a} & \color{grisclaro}b & \color{grisclaro}b & \textbf{\color{azuldos}a} & \color{grisclaro}b & \textbf{\color{azuldos}a} \\
Complejidad & \color{grisclaro}b & \color{grisclaro}b & \color{grisclaro}b & \color{grisclaro}b & \color{grisclaro}b & \color{grisclaro}b & \color{grisclaro}b & a & \color{grisclaro}b & \color{grisclaro}b & \color{grisclaro}b & \color{grisclaro}b & \color{grisclaro}b & \color{grisclaro}b & \color{grisclaro}b & \textbf{\color{azuldos}a} & \color{grisclaro}b & \color{grisclaro}b & \textbf{\color{azuldos}a} & \color{grisclaro}b & \textbf{\color{azuldos}a} \\
\#CCrea & \color{grisclaro}b & \color{grisclaro}b & \color{grisclaro}b & \color{grisclaro}b & \color{grisclaro}b & \color{grisclaro}b & \color{grisclaro}b & a & \color{grisclaro}b & \color{grisclaro}b & \color{grisclaro}b & \color{grisclaro}b & \color{grisclaro}b & \color{grisclaro}b & \color{grisclaro}b & \textbf{\color{azuldos}a} & \color{grisclaro}b & \color{grisclaro}b & \textbf{\color{azuldos}a} & \color{grisclaro}b & \textbf{\color{azuldos}a} \\
\#S & \color{grisclaro}b & \color{grisclaro}b & \color{grisclaro}b & \color{grisclaro}b & \color{grisclaro}b & \color{grisclaro}b & \color{grisclaro}b & \color{grisclaro}b & a & a & \textbf{\color{azuldos}a} & \textbf{\color{azuldos}a} & \textbf{\color{azuldos}a} & \color{grisclaro}b & \color{grisclaro}b & \textbf{\color{azuldos}a} & \textbf{\color{azuldos}a} & \color{grisclaro}b & \textbf{\color{azuldos}a} & \color{grisclaro}b & \textbf{\color{azuldos}a} \\
FInicio & \color{grisclaro}b & \color{grisclaro}b & \color{grisclaro}b & \color{grisclaro}b & \color{grisclaro}b & \color{grisclaro}b & \color{grisclaro}b & \color{grisclaro}b & a & \textbf{\color{azuldos}a} & \textbf{\color{azuldos}a} & \textbf{\color{azuldos}a} & \textbf{\color{azuldos}a} & \color{grisclaro}b & \color{grisclaro}b & \textbf{\color{azuldos}a} & \textbf{\color{azuldos}a} & \color{grisclaro}b & \textbf{\color{azuldos}a} & \color{grisclaro}b & \textbf{\color{azuldos}a} \\
FFin & \color{grisclaro}b & \color{grisclaro}b & \color{grisclaro}b & \color{grisclaro}b & \color{grisclaro}b & \color{grisclaro}b & \color{grisclaro}b & \color{grisclaro}b & a & \textbf{\color{azuldos}a} & \textbf{\color{azuldos}a} & \textbf{\color{azuldos}a} & \textbf{\color{azuldos}a} & \color{grisclaro}b & \color{grisclaro}b & \textbf{\color{azuldos}a} & \textbf{\color{azuldos}a} & \color{grisclaro}b & \textbf{\color{azuldos}a} & \color{grisclaro}b & \textbf{\color{azuldos}a} \\
\#H & \color{grisclaro}b & \color{grisclaro}b & \color{grisclaro}b & \color{grisclaro}b & \color{grisclaro}b & \color{grisclaro}b & \color{grisclaro}b & \color{grisclaro}b & \color{grisclaro}b & a & a & a & a & \color{grisclaro}b & \color{grisclaro}b & a & a & \color{grisclaro}b & a & \color{grisclaro}b & a \\
FCreacion & \color{grisclaro}b & \color{grisclaro}b & \color{grisclaro}b & \color{grisclaro}b & \color{grisclaro}b & \color{grisclaro}b & \color{grisclaro}b & \color{grisclaro}b & \color{grisclaro}b & a & \textbf{\color{azuldos}a} & \textbf{\color{azuldos}a} & \textbf{\color{azuldos}a} & \color{grisclaro}b & \color{grisclaro}b & \textbf{\color{azuldos}a} & \textbf{\color{azuldos}a} & \color{grisclaro}b & \textbf{\color{azuldos}a} & \color{grisclaro}b & \textbf{\color{azuldos}a} \\
PTeorica & \color{grisclaro}b & \color{grisclaro}b & \color{grisclaro}b & \color{grisclaro}b & \color{grisclaro}b & \color{grisclaro}b & \color{grisclaro}b & \color{grisclaro}b & \color{grisclaro}b & a & \textbf{\color{azuldos}a} & \textbf{\color{azuldos}a} & \textbf{\color{azuldos}a} & \color{grisclaro}b & \color{grisclaro}b & \textbf{\color{azuldos}a} & \textbf{\color{azuldos}a} & \color{grisclaro}b & \textbf{\color{azuldos}a} & \color{grisclaro}b & \textbf{\color{azuldos}a} \\
PPractica & \color{grisclaro}b & \color{grisclaro}b & \color{grisclaro}b & \color{grisclaro}b & \color{grisclaro}b & \color{grisclaro}b & \color{grisclaro}b & \color{grisclaro}b & \color{grisclaro}b & a & \textbf{\color{azuldos}a} & \textbf{\color{azuldos}a} & \textbf{\color{azuldos}a} & \color{grisclaro}b & \color{grisclaro}b & \textbf{\color{azuldos}a} & \textbf{\color{azuldos}a} & \color{grisclaro}b & \textbf{\color{azuldos}a} & \color{grisclaro}b & \textbf{\color{azuldos}a} \\
PLaboratorio & \color{grisclaro}b & \color{grisclaro}b & \color{grisclaro}b & \color{grisclaro}b & \color{grisclaro}b & \color{grisclaro}b & \color{grisclaro}b & \color{grisclaro}b & \color{grisclaro}b & a & \textbf{\color{azuldos}a} & \textbf{\color{azuldos}a} & \textbf{\color{azuldos}a} & \color{grisclaro}b & \color{grisclaro}b & \textbf{\color{azuldos}a} & \textbf{\color{azuldos}a} & \color{grisclaro}b & \textbf{\color{azuldos}a} & \color{grisclaro}b & \textbf{\color{azuldos}a} \\
Cobertura & \color{grisclaro}b & \color{grisclaro}b & \color{grisclaro}b & \color{grisclaro}b & \color{grisclaro}b & \color{grisclaro}b & \color{grisclaro}b & \color{grisclaro}b & \color{grisclaro}b & a & \textbf{\color{azuldos}a} & \textbf{\color{azuldos}a} & \textbf{\color{azuldos}a} & \color{grisclaro}b & \color{grisclaro}b & \textbf{\color{azuldos}a} & \textbf{\color{azuldos}a} & \color{grisclaro}b & \textbf{\color{azuldos}a} & \color{grisclaro}b & \textbf{\color{azuldos}a} \\
Cantidad & \color{grisclaro}b & \color{grisclaro}b & \color{grisclaro}b & \color{grisclaro}b & \color{grisclaro}b & \color{grisclaro}b & \color{grisclaro}b & \color{grisclaro}b & \color{grisclaro}b & a & \textbf{\color{azuldos}a} & \textbf{\color{azuldos}a} & \textbf{\color{azuldos}a} & \color{grisclaro}b & \color{grisclaro}b & \textbf{\color{azuldos}a} & \textbf{\color{azuldos}a} & \color{grisclaro}b & \textbf{\color{azuldos}a} & \color{grisclaro}b & \textbf{\color{azuldos}a} \\
\#CGen & \color{grisclaro}b & \color{grisclaro}b & \color{grisclaro}b & \color{grisclaro}b & \color{grisclaro}b & \color{grisclaro}b & \color{grisclaro}b & \color{grisclaro}b & \color{grisclaro}b & a & \textbf{\color{azuldos}a} & \textbf{\color{azuldos}a} & \textbf{\color{azuldos}a} & \color{grisclaro}b & \color{grisclaro}b & \textbf{\color{azuldos}a} & \textbf{\color{azuldos}a} & \color{grisclaro}b & \textbf{\color{azuldos}a} & \color{grisclaro}b & \textbf{\color{azuldos}a} \\
\#HAnt & \color{grisclaro}b & \color{grisclaro}b & \color{grisclaro}b & \color{grisclaro}b & \color{grisclaro}b & \color{grisclaro}b & \color{grisclaro}b & \color{grisclaro}b & \color{grisclaro}b & \textbf{\color{azuldos}a} & \textbf{\color{azuldos}a} & \textbf{\color{azuldos}a} & a & \color{grisclaro}b & \color{grisclaro}b & \textbf{\color{azuldos}a} & \textbf{\color{azuldos}a} & \color{grisclaro}b & \textbf{\color{azuldos}a} & \color{grisclaro}b & \textbf{\color{azuldos}a} \\
FValidacion & \color{grisclaro}b & \color{grisclaro}b & \color{grisclaro}b & \color{grisclaro}b & \color{grisclaro}b & \color{grisclaro}b & \color{grisclaro}b & \color{grisclaro}b & \color{grisclaro}b & \color{grisclaro}b & a & \color{grisclaro}b & \color{grisclaro}b & \color{grisclaro}b & \color{grisclaro}b & \color{grisclaro}b & \color{grisclaro}b & \color{grisclaro}b & \color{grisclaro}b & \color{grisclaro}b & \textbf{\color{azuldos}a} \\
ObsV & \color{grisclaro}b & \color{grisclaro}b & \color{grisclaro}b & \color{grisclaro}b & \color{grisclaro}b & \color{grisclaro}b & \color{grisclaro}b & \color{grisclaro}b & \color{grisclaro}b & \color{grisclaro}b & a & \color{grisclaro}b & \color{grisclaro}b & \color{grisclaro}b & \color{grisclaro}b & \color{grisclaro}b & \color{grisclaro}b & \color{grisclaro}b & \color{grisclaro}b & \color{grisclaro}b & \textbf{\color{azuldos}a} \\
FRevision & \color{grisclaro}b & \color{grisclaro}b & \color{grisclaro}b & \color{grisclaro}b & \color{grisclaro}b & \color{grisclaro}b & \color{grisclaro}b & \color{grisclaro}b & \color{grisclaro}b & \textbf{\color{azuldos}a} & \textbf{\color{azuldos}a} & a & \textbf{\color{azuldos}a} & \color{grisclaro}b & \color{grisclaro}b & \textbf{\color{azuldos}a} & \textbf{\color{azuldos}a} & \color{grisclaro}b & \textbf{\color{azuldos}a} & \color{grisclaro}b & \textbf{\color{azuldos}a} \\
Decision & \color{grisclaro}b & \color{grisclaro}b & \color{grisclaro}b & \color{grisclaro}b & \color{grisclaro}b & \color{grisclaro}b & \color{grisclaro}b & \color{grisclaro}b & \color{grisclaro}b & \textbf{\color{azuldos}a} & \textbf{\color{azuldos}a} & a & \textbf{\color{azuldos}a} & \color{grisclaro}b & \color{grisclaro}b & \textbf{\color{azuldos}a} & \textbf{\color{azuldos}a} & \color{grisclaro}b & \textbf{\color{azuldos}a} & \color{grisclaro}b & \textbf{\color{azuldos}a} \\
ObsR & \color{grisclaro}b & \color{grisclaro}b & \color{grisclaro}b & \color{grisclaro}b & \color{grisclaro}b & \color{grisclaro}b & \color{grisclaro}b & \color{grisclaro}b & \color{grisclaro}b & \textbf{\color{azuldos}a} & \textbf{\color{azuldos}a} & a & \textbf{\color{azuldos}a} & \color{grisclaro}b & \color{grisclaro}b & \textbf{\color{azuldos}a} & \textbf{\color{azuldos}a} & \color{grisclaro}b & \textbf{\color{azuldos}a} & \color{grisclaro}b & \textbf{\color{azuldos}a} \\
\#Au & \color{grisclaro}b & \color{grisclaro}b & \color{grisclaro}b & \color{grisclaro}b & \color{grisclaro}b & \color{grisclaro}b & \color{grisclaro}b & \color{grisclaro}b & \color{grisclaro}b & \color{grisclaro}b & \color{grisclaro}b & \color{grisclaro}b & \color{grisclaro}b & a & \color{grisclaro}b & a & \color{grisclaro}b & \color{grisclaro}b & \textbf{\color{azuldos}a} & \color{grisclaro}b & \textbf{\color{azuldos}a} \\
Edificio & \color{grisclaro}b & \color{grisclaro}b & \color{grisclaro}b & \color{grisclaro}b & \color{grisclaro}b & \color{grisclaro}b & \color{grisclaro}b & \color{grisclaro}b & \color{grisclaro}b & \color{grisclaro}b & \color{grisclaro}b & \color{grisclaro}b & \color{grisclaro}b & a & \color{grisclaro}b & \textbf{\color{azuldos}a} & \color{grisclaro}b & \color{grisclaro}b & \textbf{\color{azuldos}a} & \color{grisclaro}b & \textbf{\color{azuldos}a} \\
Capacidad & \color{grisclaro}b & \color{grisclaro}b & \color{grisclaro}b & \color{grisclaro}b & \color{grisclaro}b & \color{grisclaro}b & \color{grisclaro}b & \color{grisclaro}b & \color{grisclaro}b & \color{grisclaro}b & \color{grisclaro}b & \color{grisclaro}b & \color{grisclaro}b & a & \color{grisclaro}b & \textbf{\color{azuldos}a} & \color{grisclaro}b & \color{grisclaro}b & \textbf{\color{azuldos}a} & \color{grisclaro}b & \textbf{\color{azuldos}a} \\
\#F & \color{grisclaro}b & \color{grisclaro}b & \color{grisclaro}b & \color{grisclaro}b & \color{grisclaro}b & \color{grisclaro}b & \color{grisclaro}b & \color{grisclaro}b & \color{grisclaro}b & \color{grisclaro}b & \color{grisclaro}b & \color{grisclaro}b & \color{grisclaro}b & \color{grisclaro}b & a & a & \color{grisclaro}b & \color{grisclaro}b & \textbf{\color{azuldos}a} & \color{grisclaro}b & \textbf{\color{azuldos}a} \\
Dia & \color{grisclaro}b & \color{grisclaro}b & \color{grisclaro}b & \color{grisclaro}b & \color{grisclaro}b & \color{grisclaro}b & \color{grisclaro}b & \color{grisclaro}b & \color{grisclaro}b & \color{grisclaro}b & \color{grisclaro}b & \color{grisclaro}b & \color{grisclaro}b & \color{grisclaro}b & a & \textbf{\color{azuldos}a} & \color{grisclaro}b & \color{grisclaro}b & \textbf{\color{azuldos}a} & \color{grisclaro}b & \textbf{\color{azuldos}a} \\
HInicio & \color{grisclaro}b & \color{grisclaro}b & \color{grisclaro}b & \color{grisclaro}b & \color{grisclaro}b & \color{grisclaro}b & \color{grisclaro}b & \color{grisclaro}b & \color{grisclaro}b & \color{grisclaro}b & \color{grisclaro}b & \color{grisclaro}b & \color{grisclaro}b & \color{grisclaro}b & a & \textbf{\color{azuldos}a} & \color{grisclaro}b & \color{grisclaro}b & \textbf{\color{azuldos}a} & \color{grisclaro}b & \textbf{\color{azuldos}a} \\
HFin & \color{grisclaro}b & \color{grisclaro}b & \color{grisclaro}b & \color{grisclaro}b & \color{grisclaro}b & \color{grisclaro}b & \color{grisclaro}b & \color{grisclaro}b & \color{grisclaro}b & \color{grisclaro}b & \color{grisclaro}b & \color{grisclaro}b & \color{grisclaro}b & \color{grisclaro}b & a & \textbf{\color{azuldos}a} & \color{grisclaro}b & \color{grisclaro}b & \textbf{\color{azuldos}a} & \color{grisclaro}b & \textbf{\color{azuldos}a} \\
Fecha & \color{grisclaro}b & \color{grisclaro}b & \color{grisclaro}b & \color{grisclaro}b & \color{grisclaro}b & \color{grisclaro}b & \color{grisclaro}b & \color{grisclaro}b & \color{grisclaro}b & \color{grisclaro}b & \color{grisclaro}b & \color{grisclaro}b & \color{grisclaro}b & \color{grisclaro}b & \color{grisclaro}b & \color{grisclaro}b & \color{grisclaro}b & a & a & \color{grisclaro}b & a \\
EstadoA & \color{grisclaro}b & \color{grisclaro}b & \color{grisclaro}b & \color{grisclaro}b & \color{grisclaro}b & \color{grisclaro}b & \color{grisclaro}b & \color{grisclaro}b & \color{grisclaro}b & \color{grisclaro}b & \color{grisclaro}b & \color{grisclaro}b & \color{grisclaro}b & \color{grisclaro}b & \color{grisclaro}b & \color{grisclaro}b & \color{grisclaro}b & \color{grisclaro}b & a & \color{grisclaro}b & \textbf{\color{azuldos}a} \\
\#Sol & \color{grisclaro}b & \color{grisclaro}b & \color{grisclaro}b & \color{grisclaro}b & \color{grisclaro}b & \color{grisclaro}b & \color{grisclaro}b & \color{grisclaro}b & \color{grisclaro}b & \color{grisclaro}b & \color{grisclaro}b & \color{grisclaro}b & \color{grisclaro}b & \color{grisclaro}b & \color{grisclaro}b & \color{grisclaro}b & \color{grisclaro}b & \color{grisclaro}b & \color{grisclaro}b & a & a \\
Motivo & \color{grisclaro}b & \color{grisclaro}b & \color{grisclaro}b & \color{grisclaro}b & \color{grisclaro}b & \color{grisclaro}b & \color{grisclaro}b & \color{grisclaro}b & \color{grisclaro}b & \color{grisclaro}b & \color{grisclaro}b & \color{grisclaro}b & \color{grisclaro}b & \color{grisclaro}b & \color{grisclaro}b & \color{grisclaro}b & \color{grisclaro}b & \color{grisclaro}b & \color{grisclaro}b & a & \textbf{\color{azuldos}a} \\
FSolicitud & \color{grisclaro}b & \color{grisclaro}b & \color{grisclaro}b & \color{grisclaro}b & \color{grisclaro}b & \color{grisclaro}b & \color{grisclaro}b & \color{grisclaro}b & \color{grisclaro}b & \color{grisclaro}b & \color{grisclaro}b & \color{grisclaro}b & \color{grisclaro}b & \color{grisclaro}b & \color{grisclaro}b & \color{grisclaro}b & \color{grisclaro}b & \color{grisclaro}b & \color{grisclaro}b & a & \textbf{\color{azuldos}a} \\
EstadoS & \color{grisclaro}b & \color{grisclaro}b & \color{grisclaro}b & \color{grisclaro}b & \color{grisclaro}b & \color{grisclaro}b & \color{grisclaro}b & \color{grisclaro}b & \color{grisclaro}b & \color{grisclaro}b & \color{grisclaro}b & \color{grisclaro}b & \color{grisclaro}b & \color{grisclaro}b & \color{grisclaro}b & \color{grisclaro}b & \color{grisclaro}b & \color{grisclaro}b & \color{grisclaro}b & a & \textbf{\color{azuldos}a} \\
FResolucion & \color{grisclaro}b & \color{grisclaro}b & \color{grisclaro}b & \color{grisclaro}b & \color{grisclaro}b & \color{grisclaro}b & \color{grisclaro}b & \color{grisclaro}b & \color{grisclaro}b & \color{grisclaro}b & \color{grisclaro}b & \color{grisclaro}b & \color{grisclaro}b & \color{grisclaro}b & \color{grisclaro}b & \color{grisclaro}b & \color{grisclaro}b & \color{grisclaro}b & \color{grisclaro}b & a & \textbf{\color{azuldos}a} \\
\#HOri & \color{grisclaro}b & \color{grisclaro}b & \color{grisclaro}b & \color{grisclaro}b & \color{grisclaro}b & \color{grisclaro}b & \color{grisclaro}b & \color{grisclaro}b & \color{grisclaro}b & \color{grisclaro}b & \color{grisclaro}b & \color{grisclaro}b & \color{grisclaro}b & \color{grisclaro}b & \color{grisclaro}b & \color{grisclaro}b & \color{grisclaro}b & \color{grisclaro}b & \color{grisclaro}b & a & \textbf{\color{azuldos}a} \\
\#HDes & \color{grisclaro}b & \color{grisclaro}b & \color{grisclaro}b & \color{grisclaro}b & \color{grisclaro}b & \color{grisclaro}b & \color{grisclaro}b & \color{grisclaro}b & \color{grisclaro}b & \color{grisclaro}b & \color{grisclaro}b & \color{grisclaro}b & \color{grisclaro}b & \color{grisclaro}b & \color{grisclaro}b & \color{grisclaro}b & \color{grisclaro}b & \color{grisclaro}b & \color{grisclaro}b & a & \textbf{\color{azuldos}a} \\
\#CAt & \color{grisclaro}b & \color{grisclaro}b & \color{grisclaro}b & \color{grisclaro}b & \color{grisclaro}b & \color{grisclaro}b & \color{grisclaro}b & \color{grisclaro}b & \color{grisclaro}b & \color{grisclaro}b & \color{grisclaro}b & \color{grisclaro}b & \color{grisclaro}b & \color{grisclaro}b & \color{grisclaro}b & \color{grisclaro}b & \color{grisclaro}b & \color{grisclaro}b & \color{grisclaro}b & a & \textbf{\color{azuldos}a} \\
\#CApo & \color{grisclaro}b & \color{grisclaro}b & \color{grisclaro}b & \color{grisclaro}b & \color{grisclaro}b & \color{grisclaro}b & \color{grisclaro}b & \color{grisclaro}b & \color{grisclaro}b & \color{grisclaro}b & \color{grisclaro}b & \color{grisclaro}b & \color{grisclaro}b & \color{grisclaro}b & \color{grisclaro}b & \color{grisclaro}b & \color{grisclaro}b & \color{grisclaro}b & \color{grisclaro}b & a & \textbf{\color{azuldos}a} \\
\end{longtable}}

